\subsection{Fresh reconstruction of the parameter-preserving PL field comparison}

\begin{saresult}{A.115R}{Proved under stated hypotheses}{Finite PL families descend across real closed fields}
Let \(F\subset K\) be an inclusion of real closed fields. Let \(N\) be a compact finite PL polyhedron defined over \(F\), let \(A\subset N\) be a subpolyhedron, and let \(M_1,\ldots,M_s\) be finitely many marked subpolyhedra, all defined over \(F\). Consider the actual simplicial groups of jointly PL, parameter-preserving automorphisms of \(N\), fixing \(A\) pointwise and preserving the \(M_i\) setwise. Scalar extension from \(F\) to \(K\) induces a weak equivalence of these simplicial groups.

More explicitly, every jointly \(K\)-PL family \(H:N_K\times P_K\to N_K\times P_K\), \(P=[0,1]^q\), preserving the projection to \(P\), fixing \(A_K\times P_K\), preserving the marks, and equal to the identity on \(N_K\times\partial P_K\), is joined by a jointly \(K\)-PL homotopy with the same properties to the scalar extension of such an \(F\)-PL family. Also, every \(K\)-PL filling of a specified \(F\)-PL boundary family admits an \(F\)-PL filling with exactly that boundary.
\end{saresult}

\begin{proof}
This is a fresh reconstruction, not a claim that the inaccessible earlier A.115 text has been recovered verbatim. We prove the finite simultaneous-triangulation, transfer, and isotopy steps directly. All PL maps below are finite and jointly PL in every displayed parameter.

Choose a finite simplicial complex \(\mathcal Q\) over \(F\) with underlying polyhedron \(N\), compatible with \(A,M_1,\ldots,M_s\). Give \(P\) its finite face complex. The fixed product cells are
\[
 D\times\sigma,\qquad D\in\mathcal Q,\quad
 \sigma\text{ a face of }P.
 \tag{115R.1}
\]
They are convex polytopes in their affine hulls. Write
\[
 E=(A\times P)\cup(N\times\partial P).
 \tag{115R.2}
\]
For \(q=0\), omit the parameter-boundary term.

\textbf{Simultaneous triangulations over one parameter complex.}
First choose a finite target polytopal subdivision on which \(H^{-1}\) is affine, refining all fixed target product cells and all their images under \(H\). Such a subdivision is obtained from the finitely many linearity polyhedra of \(H^{-1}\), the fixed cells, and their finitely many polyhedral images, by a common finite hyperplane arrangement. Include the spatial marks and \(E\). The inverse-image subdivision under \(H\) then refines the same fixed source product cells.

We can refine the target subdivision so that its projection \(p:N\times P\to P\) is simplicial over a finite triangulation \(\mathcal L\) of \(P\). Here is a finite construction. Subdivide \(P\) by the affine hulls and extended supporting hyperplanes of all projections of target cells and their faces. Thus every projected polytope is a union of cells in a common base polytopal complex. Cut each target cell by the inverse images of these base cells. Every nonempty resulting cell maps onto a base cell, and its faces map onto faces of that cell. Choose one relative-interior point \(b_\sigma\) in each base cell. In every target cell \(C\) projecting onto \(\sigma\), choose a relative-interior point \(c_C\) with \(p(c_C)=b_\sigma\). This is possible because a linear map sends the relative interior of a convex polytope onto the relative interior of its image; this follows directly by positive convex combinations of vertices. Cone the faces successively from these points. A chain of target cells maps to a chain of base cells, with repetitions allowed, so the resulting triangulations make \(p\) simplicial. This construction is finite over any real closed field.

Pull the resulting target triangulation back through \(H\). Because \(H^{-1}\) is affine on its cells, this gives a source triangulation \(\mathcal S\), a target triangulation \(\mathcal T\), and a common base triangulation \(\mathcal L\), for which
\[
 H:\mathcal S\longrightarrow\mathcal T,\qquad
 p:\mathcal S\longrightarrow\mathcal L,\qquad
 p:\mathcal T\longrightarrow\mathcal L
 \tag{115R.3}
\]
are simplicial, and the first map is a simplicial isomorphism. Both source and target subdivisions refine (115R.1). On \(E\), where \(H=\operatorname{id}\), their triangulations agree literally. All marked subpolyhedra are subcomplexes.

\textbf{Transferring the finite realization data.}
Retain the finite abstract complexes, their vertex maps in (115R.3), the carriers in the fixed product cells, and the identifications of the source and target vertices on \(E\). Their geometric realization over the fixed \(F\)-polyhedron is a first-order condition in the coordinates of finitely many vertices. It says that each abstract simplex is affinely independent, that pairwise intersections are exactly the prescribed common faces, that the complexes cover \(N\times P\) and \(P\), that their assigned subcomplexes cover each prescribed fixed cell and mark, and that the projection of each vertex is the prescribed base vertex. These conditions use only ordered-field formulas over \(F\). For example, membership in a simplex is an existential barycentric-coordinate formula, and equality of its intersection or union with another specified polyhedron is expressed by a universal point quantifier. Equality of the source and target coordinates on \(E\) is included explicitly.

The original \(K\)-coordinates satisfy this finite formula. The elementary inclusion of real closed fields \(F\subset K\) therefore gives a realization with coordinates in \(F\). Denote the transferred complexes by
\(\mathcal S_F,\mathcal T_F,\mathcal L_F\), and their transferred simplicial isomorphism by \(H_F\). It is a jointly \(F\)-PL family over the parameter projection and has all prescribed fixed and marked conditions.

Let \(R_s,R_t\) be the simplicial isomorphisms from the original source and target realizations to the transferred ones, viewed over \(K\); let \(\beta\) be the analogous base isomorphism. They preserve every fixed product cell, and \(\beta\) preserves every face of \(P\). We have exact identities
\[
 pR_s=\beta p,\qquad pR_t=\beta p,\qquad
 R_tH=H_FR_s,\qquad R_s|_E=R_t|_E.
 \tag{115R.4}
\]
The first two identities hold on every simplex, not just at its vertices: both sides are affine on that simplex, since the projection and the base map are simplicial, and they agree on its vertices. The last identity uses the literal agreement of both subdivisions and the transferred coordinates on \(E\). Neither \(R_s\) nor \(R_t\) is asserted to be the identity there.

\textbf{A finite fibrewise cell contraction.}
We use the following elementary fact. If a jointly PL homeomorphism \(a:N\times P\to N\times P\) preserves the projection and every product cell (115R.1), then it is joined to the identity by a jointly PL homotopy through such maps. If \(a\) fixes a union \(E_0\) of product cells pointwise, the homotopy can fix \(E_0\) pointwise.

Construct it by increasing total dimension of \(D\times\sigma\). Prescribe the initial map \(a\) at time zero, the identity at time one, and all homotopies already assigned on proper faces. On cells contained in \(E_0\), prescribe the identity throughout. For a remaining product cell, the prescribed data form a PL homeomorphism of the entire boundary of
\[
 D\times\sigma\times[0,1].
\]
They preserve the projection to \(\sigma\times[0,1]\). Cone this TOTAL boundary homeomorphism from the source center
\((d_0,b_\sigma,\tfrac12)\) to the identical target center, after triangulating the boundary. This extends to a PL homeomorphism of the full product polytope. It preserves the parameter and time exactly: affine projection takes the source and target cone segments to the same segment with endpoints \((b_\sigma,\tfrac12)\) and the boundary parameter. Thus it is a jointly PL family, not a formula obtained by multiplying a parameterized spatial map by a varying scalar. Finite induction and pasting give the required homotopy. The same construction, with only the parameter cells, joins a face-preserving PL homeomorphism of \(P\) to the identity through face-preserving PL maps.

\textbf{Untwisting the transferred family.}
Put \(B=\operatorname{id}_N\times\beta\), and define
\[
 a=B^{-1}R_s,\qquad k=R_tR_s^{-1}.
\]
Both preserve the projection to \(P\) and all fixed product cells. The map \(k\) fixes \(E\) pointwise by (115R.4). We have
\[
 H_F=k\,B\,a\,H\,a^{-1}B^{-1}.
 \tag{115R.5}
\]
Contract \(k\) to the identity through cell-preserving fibrewise PL maps fixed on \(E\), and contract \(a\) to the identity through cell-preserving fibrewise PL maps. Along the simultaneous substitution in (115R.5), the quotient remains the identity on \(E\): the first contraction fixes \(E\), while conjugation by the second preserves \(E\) setwise. Every spatial mark is preserved. The resulting endpoint is \(BHB^{-1}\).

Finally choose the face-preserving PL isotopy \(\beta_t\) from \(\beta\) to the identity supplied by the preceding cell argument, and put \(B_t=\operatorname{id}_N\times\beta_t\). The families \(B_tHB_t^{-1}\) preserve the original parameter projection and all based/fixed conditions, and end at \(H\). They are jointly PL: the total map \((x,p,t)\mapsto(x,\beta_t(p),t)\) is PL, and inverse and composition of finite PL homeomorphisms are PL. Concatenating and reversing these homotopies joins \(H\) to the scalar extension of \(H_F\) with all required conditions. This proves surjectivity on components and on all homotopy groups represented by based cubes.

\textbf{Exact descent of fillings.}
Let a prescribed boundary family \(G\) be jointly PL over \(F\), and suppose a jointly \(K\)-PL filling exists. Fix the actual finite format of that filling: a finite union of closed simplices describing its graph, or equivalently finitely many affine pieces with finitely many vertex coordinates. Treat those finitely many \(K\)-coefficients as variables. The assertion that they define a continuous graph whose projections to source and target are bijective, preserve the parameter coordinates, fix \(A\), preserve the marks, and equal the given semilinear graph of \(G\) on the parameter boundary is first-order over \(F\). The source and target are compact, so the two continuous bijective graph projections give a homeomorphism; alternatively continuity and inverse continuity may both be included directly in their epsilon-delta first-order formulas. The graph is a finite union of simplices, so these bijective projections define a PL homeomorphism: each simplex projection is injective, its inverse is affine on its image, and a common finite polyhedral refinement gives the PL map and inverse.

This existential coefficient formula is true in \(K\), hence in \(F\). Its \(F\)-solution is the desired exact \(F\)-PL filling. In particular, a based \(F\)-PL cube nullhomotopic over \(K\) is nullhomotopic over \(F\), and the same argument applies to homotopies between two \(F\)-PL vertices. The actual simplicial PL groups are Kan groups; the usual based-cube description of their homotopy groups therefore proves the stated weak equivalence. Empty \(N\) gives the constant one-element simplicial group; zero-dimensional parameters and disconnected polyhedra are already included in the finite construction.
\end{proof}

\paragraph{Consequence needed for the all-field relative comparison.}
On a rational PL model, take \(F\) to be the real algebraic field inside \(K\). Every based PL homotopy class then has a representative whose finitely many coefficients lie in \(F\). Its compact Lipschitz bounds are bounded by standard integers. This conclusion uses the full parameter-preserving finite construction above; it is not inferred from an ordinary, unparameterized PL isotopy theorem.
